
\begin{theorem}[primitive--Dirichlet 级数的自然边界：\texorpdfstring{$\Re(s)=0$}{Re(s)=0}]\label{thm:real-input-40-primedirichlet-dense-branch}
令
\[
\mathcal{P}(s):=P(\lambda^{-s})=\sum_{n\ge 1}p_n\,\lambda^{-ns}\qquad(\Re(s)>1),
\]
其中 $P(z)=\sum_{n\ge 1}p_n z^n$ 为 primitive 生成函数，$\lambda=\rho(M)=\varphi^2$，并定义
\[
\widehat{\zeta}(s):=\zeta_M(\lambda^{-s}).
\]
则：
\begin{enumerate}
  \item \textbf{（解析延拓到右半平面）} 对任意 $\sigma_0>0$，级数
  \[
  \sum_{k\ge 1}\frac{\mu(k)}{k}\log \widehat{\zeta}(ks)
  \]
  在集合
  \(
  \{s:\Re(s)\ge \sigma_0\}
  \)
  的任意紧子集上（避开下述极点列）一致收敛，从而定义出 $\mathcal{P}(s)$ 在开半平面 $\Re(s)>0$ 上的一个解析延拓（其在 $\Re(s)>1$ 与原 Dirichlet 级数一致）。
  \item \textbf{（稠密奇点列与自然边界）} 对每个奇素数 $p$ 与任意 $m\in\ZZ$，点
  \[
  s_{p,m}:=\frac{1}{p}+\frac{2\pi i m}{p\log\lambda}
  \]
  是该解析延拓的一个对数型分支奇点。并且集合 $\{s_{p,m}\}_{p\ \mathrm{odd\ prime},\,m\in\ZZ}$ 在直线 $\Re(s)=0$ 上稠密。
\end{enumerate}
因此 \(\Re(s)=0\) 是 \(\mathcal{P}(s)\) 的\textbf{自然边界}：\(\mathcal{P}(s)\) 不可能在任何虚轴点的邻域内作解析延拓穿过该直线。
\end{theorem}

\begin{proof}
\emph{第 1 点（收敛与解析性）}：
当 $|z|$ 足够小时有 $\log\zeta_M(z)=O(z)$（因为 $\zeta_M(0)=1$ 且 $\zeta_M$ 在 $0$ 的邻域解析），因此对任意 $\sigma_0>0$，当 $\Re(s)\ge \sigma_0$ 时
\(
|\log\widehat{\zeta}(ks)|=|\log\zeta_M(\lambda^{-ks})|=O(\lambda^{-k\sigma_0})
\)
随 $k$ 指数衰减。于是 $\sum_{k\ge 1}|\mu(k)|\,|\log\widehat{\zeta}(ks)|/k$ 在任意避开极点列的紧集上一致收敛；逐项可导与一致收敛给出其和函数在 $\Re(s)>0$（去掉极点列）上解析，从而得到 $\mathcal{P}(s)$ 的解析延拓。

\emph{第 2 点（显式奇点与稠密性）}：
由本核的行列式分解
\[
\det(I-zM)=(1-z)^2(1+z)^3(1-3z+z^2)(1+z-z^2),
\]
可见 $z_\star=\lambda^{-1}$ 是 $\zeta_M(z)$ 的一个\textbf{简单极点}（对应因子 $1-3z+z^2$ 的简单零点），故 $\log\zeta_M(z)$ 在 $z=z_\star$ 具有对数型分支奇性。
当 $s=s_{p,m}$ 时有
\[
\lambda^{-ps}=\lambda^{-1}e^{-2\pi i m}=z_\star,
\]
因此 $\log\widehat{\zeta}(ps)=\log\zeta_M(\lambda^{-ps})$ 在 $s=s_{p,m}$ 处为对数型分支奇点，从而项
\(
\frac{\mu(p)}{p}\log\widehat{\zeta}(ps)
\)
在该点奇异。对任意 $k\neq p$，若 $\log\widehat{\zeta}(ks)$ 在 $s=s_{p,m}$ 奇异，则必有 $\lambda^{-ks_{p,m}}=z_\star$，等价于 $ks_{p,m}=1+\frac{2\pi i \ell}{\log\lambda}$（某个 $\ell\in\ZZ$）。比较实部得 $k/p=1$，故 $k=p$。因此除 $k=p$ 一项外，其余各项在 $s_{p,m}$ 的某邻域解析，且由第 1 点一致收敛可知其和亦解析，于是 $\mathcal{P}(s)$ 在 $s_{p,m}$ 必为对数型分支奇点。

最后证明稠密性：给定任意 $t\in\RR$ 与 $\varepsilon>0$，取足够大的奇素数 $p$，令
\(
m:=\left\lfloor \frac{t\,p\log\lambda}{2\pi}\right\rceil
\)
为最近整数，则
\[
\left|\Im(s_{p,m})-t\right|
=\left|\frac{2\pi m}{p\log\lambda}-t\right|
\le \frac{\pi}{p\log\lambda}<\varepsilon
\]
当 $p$ 足够大时成立，同时 $\Re(s_{p,m})=1/p<\varepsilon$。故 $s_{p,m}\to it$，从而 $\{s_{p,m}\}$ 在虚轴上稠密。由于任意穿过虚轴的邻域都包含这些奇点，故无法跨越解析延拓，\(\Re(s)=0\) 为自然边界。
\end{proof}

\begin{proposition}[分支点继承（prime-zeta 缩放律）]\label{prop:primedirichlet-branch-inherit}
保持定理 \ref{thm:real-input-40-primedirichlet-dense-branch} 的记号，并假设 \(\Re(s_0)>0\)。
若 \(\widehat{\zeta}(s)\) 在 \(s=s_0\) 具有极点（或等价地 \(\log\widehat{\zeta}(s)\) 在 \(s=s_0\) 具有对数型分支奇性），则对任意奇素数 \(p\)，函数 \(\mathcal{P}(s)\) 在
\[
s=\frac{s_0}{p}
\]
处具有对数型分支奇点。
\end{proposition}
\begin{proof}
由 prime-zeta 形式恒等式
\(
\mathcal{P}(s)=\sum_{k\ge 1}\frac{\mu(k)}{k}\log\widehat{\zeta}(ks)
\)
可知：在点 \(s=s_0/p\) 处，项 \(\frac{\mu(p)}{p}\log\widehat{\zeta}(ps)\) 由于 \(ps=s_0\) 必为对数型分支奇点。
对任意 \(k\neq p\)，若 \(\log\widehat{\zeta}(ks)\) 在 \(s=s_0/p\) 奇异，则 \(ks\) 必落在 \(\widehat{\zeta}\) 的极点列上，特别满足 \(\Re(ks)=\Re(s_0)\)。
比较实部得 \(k\Re(s_0)/p=\Re(s_0)\)，从而 \(k=p\)，矛盾。
因此除 \(k=p\) 一项外，其余各项在 \(s_0/p\) 的某邻域解析；由定理 \ref{thm:real-input-40-primedirichlet-dense-branch} 第 1 点的一致收敛，解析项之和仍解析，故 \(\mathcal{P}(s)\) 在 \(s=s_0/p\) 必为对数型分支奇点。证毕。
\end{proof}

\begin{corollary}[toy-RH 临界线极点诱导的 \texorpdfstring{$\Re(s)=\frac{1}{2p}$}{Re(s)=1/(2p)} 分支点族]\label{cor:real-input-40-primedirichlet-branch-halfp}
保持定理 \ref{thm:real-input-40-primedirichlet-dense-branch} 的记号。对每个奇素数 \(p\) 与任意 \(m\in\ZZ\)，点
\[
